\subsection{The fixed-support classification}

For $S_B=\{p\text{ prime}:p\le B\}$, with $B=7,13,19$, we solve a
two-term $S$-unit problem, reconstruct the primitive triangles and their
normalized square progressions, and test every centre interaction.
The result is a complete classification for the specified supports,
with unrestricted exponents.

\begin{theorem}[Smooth differences and a complete $(9,7)$ atlas]\label{main:smoothsupport}
\textbf{Computer-assisted, using imported complete $S$-unit totals.}
Let $S_B$ be the primes at most $B$. \emph{(i)} There is no fully magic
$3\times3$ square of nine distinct positive integer squares whose
transversal difference $d$ satisfies $\supp(d)\subseteq S_{19}$.
\emph{(ii)} The $(9,7)$ arrays with $\supp(|D|)\subseteq S_B$, modulo
$D_4$ and common rational-square scaling, have exactly $3,45,384$ classes
for $B=7,13,19$, respectively, represented by primitive positive roots.
\emph{(iii)} For $B=19$, the sharp lower bound for
$\eta=|A+P-2E|/|D|$ is $347984603/9837828000$, attained by exactly one
class, already present for $B=13$. These are unbounded-exponent statements:
completeness uses Theorem~A of von K\"anel--Matschke, not a search cutoff.
\end{theorem}

\subsection{The imported total and its normalization}

Let $\mathcal A(S)$ be the set of integer triples
\[
0<a\le b<c,\qquad a+b=c,\qquad \gcd(a,b)=1,
\qquad \supp(abc)\subseteq S.
\]
Theorem~A of von K\"anel and Matschke~\cite{VKM16}, printed p.~7,
counts solutions of $x+y=1$ in rational $S$-units modulo the six usual
symmetries. Each orbit has exactly one representative $0<x\le1/2$;
writing it as $x=a/c$, $1-x=b/c$ gives exactly one element of
$\mathcal A(S)$. The exceptional orbit represented by $(1,1,2)$ is counted
once. Thus their entries for the first four, six and eight primes give
\[
|\mathcal A(S_7)|=63,\qquad |\mathcal A(S_{13})|=545,\qquad
|\mathcal A(S_{19})|=3649.
\]
The first two totals were already obtained by de Weger, as credited in
\cite{VKM16}. We import only these three totals, not a conjecture comparing
different supports of the same size.

The supplied producer finds $3649$ distinct valid triples with
$c\le10^{10}$. The separate verifier checks integrality, positivity,
ordering, coprimality, addition and prime support for every row, and checks
the imported total. A subset of a finite set with the same cardinality is
the entire set: this is the completeness argument, with \emph{no remaining
height restriction}. The production bound alone would prove no such claim.
The witnesses are regenerated locally, not copied from the authors' data
files; their source theorem remains a mathematical dependency.

\subsection{From additive triples to square progressions}

\begin{lemma}[Exact triangle reconstruction]\label{lem:trianglecatalogue}
Primitive positive integer right triangles, up to swapping the legs, whose
area is supported on $S$ correspond bijectively to pairs $m>k>0$ with
$\gcd(m,k)=1$, $m+k$ odd, and
\[
(\min(k,m-k),\max(k,m-k),m)\in\mathcal A(S),\qquad
(k,m,m+k)\in\mathcal A(S).
\]
Their legs, hypotenuse and area are
\[
\{a,b\}=\{2mk,m^2-k^2\},\qquad c=m^2+k^2,
\qquad A=mk(m-k)(m+k).
\]
\end{lemma}
\begin{proof}
Euclid's parametrization is unique with these coprimality and parity
conditions. Its area is supported on $S$ exactly when all four positive
factors $m,k,m-k,m+k$ are supported on $S$. The two displayed additive
triples express precisely these conditions; their coprimality follows
from $\gcd(m,k)=1$. Conversely, membership and odd $m+k$ give Euclid's
conditions and the required area. This proves both directions, not just
a procedure for generating examples.
\end{proof}

Write $A=n h^2$ with $n$ squarefree. The corresponding normalized square
progression has positive rational roots
\[
\left(\frac{|a-b|}{h},\frac c h,\frac{a+b}{h}\right),
\qquad d=4n,\qquad \sigma=\frac{c^2}{h^2}>4n.
\]
Conversely, divide any positive integer-root square progression by the gcd
$g$ of its roots. Write the divided roots as $r',s',t'$; they are odd,
and $((t'-r')/2,(t'+r')/2,s')$ is a primitive right triangle of area
$d/(4g^2)$.
Thus $d/4=n(hg)^2$. Progressions of the same difference necessarily have
the same $n$ and common scale $\lambda=hg$, so their actual centres are
$\lambda^2\sigma$. Consequently the full-magic interaction is exactly
\[
\sigma_1+\sigma_2=2\sigma_0.
\]
All denominators of the normalized roots divide $h$ and are supported on
$S$. Clearing them and dividing by the common root gcd therefore gives
every root-primitive integral array without introducing a prime outside
$S$ into its difference.

\subsection{The complete finite tests}

\begin{proof}[Proof of Theorem~\ref{main:smoothsupport}]
The complete additive lists and Lemma~\ref{lem:trianglecatalogue} give the
following exact counts. A midpoint test checks one unordered pair of
distinct centres within the same squarefree area class and asks whether
their average is another centre in that class.
\[
\begin{array}{c|rrrrrr}
B & |\mathcal A(S_B)| & \text{triangles} & \text{classes} &
\text{pairs} & \text{hits} & (9,7)/D_4\text{ classes}\\\hline
7  & 63   & 12  & 9   & 4   & 0 & 3\\
13 & 545  & 62  & 39  & 36  & 0 & 45\\
19 & 3649 & 232 & 133 & 180 & 0 & 384
\end{array}
\]
Every number in this table is recomputed with integer and rational
arithmetic by the companion's exact checker:
\begin{center}\texttt{scripts/verify\_smooth\_support.py}.\end{center}
In particular,
the normalized centres within each class are distinct and no midpoint
test succeeds. By the preceding equivalence, no fully magic square in
the stated domain exists. This establishes part~(i).

For part~(ii), group the triangles by $n$ and select three distinct
triangles in one group. The three normalized progressions give nine
candidate entries. Reject nonpositive or repeated entries; for these
catalogues none of the $1,15,128$ unordered triples is rejected. Choose
one of the three centres as the central cell $E$, use the other two as
$A,P$ in the normal form, and set $D=4n$. The absence of midpoint hits
ensures exactly seven equal line sums. Clear root denominators, divide
by the root gcd, and take the lexicographically least positive-root array
in its eight-element $D_4$ orbit.

These operations are exhaustive by the normal form, which also holds
when all eight lines agree. For an array with exactly one failed line,
orienting that diagonal as the antidiagonal leaves precisely the choices
of interchanging $A,P$ and changing the sign of $D$. Those choices are
already identified by $D_4$; the central cell is fixed by every grid
symmetry. Therefore the three choices of $E$ are inequivalent. The unique
primitive normalization of a positive rational-root array also removes
all common rational-square scaling. The verifier reconstructs and compares
all $3,45,384$ representatives, establishing the classification.

Finally, for each representative compute the seven-line sum $M$, the
failed sum $F$, and
\[
\eta=\frac{|M-F|}{|D|}=\frac{|A+P-2E|}{|D|}.
\]
This quotient is invariant under $D_4$ and common square scaling. Exact
comparison over the complete lists gives minima
$1727/1120$ for $S_7$ and $347984603/9837828000$ for both $S_{13}$ and
$S_{19}$, each attained by one class. For the latter, a primitive root
representative (the \emph{entries of the array are its squares}) is
\[
\begin{pmatrix}
87942&140595&157690\\
186290&132558&9555\\
99645&122590&165558
\end{pmatrix}.
\]
It has $|D|=9837828000$, $M=52366885489$, $F=52714870092$, and
$M-F=-347984603$. These identities and uniqueness are checked from the
full atlas, proving part~(iii). No floating-point comparison is used.
\end{proof}

For $S_7$ the unique unordered triple, after common root-primitive
normalization, has roots $(2,58,82)$, $(46,74,94)$ and $(97,113,127)$ and common
difference $3360$. Houston~\cite{HOU16B} already displays exactly these
three progressions: his elliptic parameter has difference $210$, and the
displayed integral roots give difference $3360$. The contribution here is
their place in the complete fixed-support classification. Their three choices of the central progression account
for the three classes. The atlas is for $(9,7)$ arrays; it is separate
from the bounded $(8,7)$ census of Appendix~A. Also, $\eta$ is normalized
by the progression difference, not by the magic sum: it is not a percentage
error or a comparison with unrestricted near-miss records.

\subsection{Consequences and limits}

Every hypothetical fully magic square of distinct positive integer
squares must have a prime at least $23$ in each of its two transversal
differences. Indeed, write its normal form as
$M(E-t,E,E+t;d)$, where $M$ denotes the matrix in
Theorem~\ref{main:normalform}. A clockwise quarter-turn gives
$M(E+d,E,E-d;t)$. Both differences are nonzero, and applying part~(i) in
each orientation proves the assertion. The two large primes need not be
different. These are primes of the \emph{differences}, not of the entries.

In particular, the fifteen supports consisting of $2,3$ and two distinct
primes chosen from $S_{19}\setminus\{2,3\}$ are excluded. This does not close arbitrary
four-prime supports, bound the number of support primes, or settle the
general magic-square-of-squares problem. The general fixed-squareclass
finiteness result of Section~\ref{sec:finiteness} remains non-effective;
the present enumeration is obtained through the separate fixed-support
route and has been executed only for the declared catalogues.
